\documentclass[12pt,letterpaper]{article}
\usepackage[T1]{fontenc}
\IfFileExists{newtxmath.sty}%
  {\usepackage{newtxtext,newtxmath}}%
  {\usepackage{mathptmx}\usepackage{amssymb}}
\usepackage[letterpaper,margin=1in]{geometry}
\usepackage{amsmath,mathtools}
\usepackage{booktabs,longtable,array,enumitem}
\usepackage{microtype}
\usepackage[hidelinks]{hyperref}
\newcommand{\Gt}{\widetilde G}
\newcommand{\At}{\widetilde A}
\newcommand{\hl}{\tilde h}
\newcommand{\wt}{\tilde w}
\newcommand{\ZA}{Z_A}
\newcommand{\Of}{\mathcal O}
\DeclareMathOperator{\tr}{tr}

\title{Revision log for\\ ``On the unitary dual of $\overline{SL(2,F)}$,
$F$ $p$-adic''\\[4pt]\large Draft 1 (November 9, 2026) to Draft 2
(November 13, 2026)}
\author{Prepared by Claude (Anthropic) for D. Joyner}
\date{November 13, 2026}

\begin{document}
\maketitle
\tableofcontents

\section{Files}

\begin{itemize}[leftmargin=*]
\item \texttt{sl2bar-unitary-dual-2026.tex}: Draft~2 of the paper.
\item \texttt{revision-log.tex}: this log.
\item \texttt{scripts/kubota\_torus\_check.py}: exact-arithmetic checks
  of the cocycle computations (Section~\ref{sec:script}).
\item Inputs:
  \begin{itemize}
  \item the 1999 manuscript \texttt{sph\_corr5b.tex}, dated 7-21-99,
  \item Draft~1, in a file whose name ends in
    \texttt{metaplectic\_cover\_2026c.tex},
  \item the style guide, Draft~2 of November 11, 2026, in
    \texttt{joyner-math-paper-style-guide.md}.
  \end{itemize}
\end{itemize}

Draft~2 compiles with \texttt{pdflatex} with no errors, no undefined
references and no overfull boxes. It was compiled here with
\texttt{mathptmx}, because \texttt{newtx} is not installed in the
sandbox. The preamble loads \texttt{newtxtext} and \texttt{newtxmath}
when they are installed, so a local build will use them. It has not
been test-compiled with \texttt{newtx}.

\section{AI usage}

\begin{itemize}[leftmargin=*]
\item \textbf{Argon-Gemini.} Refereed the 1999 manuscript, suggested
  corrections, added a glossary and produced a first revised
  manuscript. A later revision attempt was not completed. (Source: the
  abstract and acknowledgments of Draft~1.)
\item \textbf{Sol-ChatGPT (OpenAI).} Wrote an independent referee
  report and prepared Draft~1, the corrected LaTeX edition. (Source:
  Draft~1.)
\item \textbf{Claude (Anthropic), this draft.} Rewrote Draft~1 to
  follow the style guide. Checked every mathematical claim of Draft~1
  independently (Section~\ref{sec:levels}). Supplied complete proofs
  where Draft~1 gave sketches. Wrote the Background section, all new
  examples, and the expanded literature review. Checked bibliographic
  data by web search. Wrote the verification script and this log.
\end{itemize}
Claude did not see the Argon-Gemini manuscript or either referee
report. Claude's knowledge of them comes only from Draft~1.

\section{Original results and results new in the revision}

Rule~6.5 asks for this separation.

\subsection{Results of the 1999 manuscript}
\begin{enumerate}[leftmargin=*]
\item Vanishing lemma for orbital integrals off the $d$th powers. Kept
  (Lemma~6.1), with the correct quotient.
\item Surjectivity of the split orbital-integral map (sketched in 1999,
  with the final step referred to Shalika germs). Kept and now fully
  proved (Theorem~6.4).
\item Description of maximal abelian subgroups and genuine
  representations of $\At$. The 1999 ``2-1 correspondence''
  $\Pi(A_n)\to\Pi(\At)_{\mathrm{gen}}$ is wrong as stated, see
  Section~\ref{subsec:removed}. Replaced by Theorem~5.6.
\item Character formula for principal series, stated with
  $\tr\chi(a)+\tr(\chi(a)^{-1})$ and with proof omitted. Corrected to
  Weyl conjugation and proved (Theorem~8.3).
\item Character identity $\tr\overline\pi(\bar f)=\tr\pi(f)$ for
  $\gcd(p,2n)=1$ and all $n$. Correct only for $n$ odd or $4\mid n$,
  with the normalization of Theorem~9.4.
\item Unitarity of principal series, complementary series for
  $0<s<1/n$, the special representation and the ``trash''
  representation at $s=1/n$. Not proved. Recorded as cited material in
  Remarks~2026.2--2026.5.
\end{enumerate}

\subsection{New in Draft 1 (Sol-ChatGPT)}
The correct center $\ZA$ and the index $d^2$. The Stone--von Neumann
parametrization by genuine characters of $\ZA$. The distinction between
Weyl conjugation and inversion. The restriction to $n$ odd or $4\mid n$
and the power map $\nu$ (Lemma~9.1). Explicit Haar normalizations and
the $d$ versus $d^2$ distinction. Removal or demotion of the unproved
unitary-dual claims.

\subsection{New in Draft 2 (Claude)}
\begin{itemize}[leftmargin=*]
\item Lemma~3.2 (cocycle values, uniqueness of the splitting of $N$,
  centrality of the preimage of $\pm I$). New, with proof.
\item Lemma~5.1 (exact commutator $(a,b)_n^{-2}$). Draft~1 had
  $(a,b)_n^{\pm2}$.
\item Lemma~5.4 (Weyl conjugation on $\At$). Draft~1 stated (b) inside
  the proof of Lemma~8.1. Parts (a), (c), (d) are new.
\item Lemma~6.3 (constant term). Draft~1 stated the formula without
  proof.
\item Theorem~6.4. Draft~1 gave a sketch. Step~1 (submersion argument,
  fiber description, explicit preimage) and Step~2 (explicit functions
  near $\pm I$ and their transforms) are written out.
\item Proposition~7.1. Draft~1 cited Moen for irreducibility. Draft~2
  proves (a) directly and (b) from Frobenius reciprocity and the
  geometric lemma, which are cited (Remark~2026.2).
\item Lemma~7.2 (spherical vectors). New.
\item Lemma~8.1 (trace of an induced representation by the kernel
  method). New. Corollary~8.2 now follows from it directly.
\item Theorem~8.3. Draft~1 gave a sketch citing [G, F, FK]. Draft~2
  proves well-definedness, the trace identity via the Weyl integration
  formula, and local integrability.
\item Section~3 (Background) and Examples 3.1, 3.3--3.7, 5.5, 6.2, 6.6,
  7.3, 8.5, 9.3. New. Examples 5.7, 5.8, 9.5 are from Draft~1, revised.
\end{itemize}

\section{Changes from Draft 1 to Draft 2}

\subsection{Style-guide conformance}
\begin{itemize}[leftmargin=*]
\item \textbf{1.1.} Twelve-point Times (\texttt{newtx}, with
  \texttt{mathptmx} fallback), US Letter, one-inch margins. Draft~1 used
  Latin Modern.
\item \textbf{1.2, 1.3.} Compiles cleanly. Two overfull boxes were
  removed by rewording.
\item \textbf{1.5.} All cross-references use \verb|\ref| and all
  citations use \verb|\cite|. Section titles with math use
  \verb|\texorpdfstring|.
\item \textbf{1.7.} Citation keys changed to author letters plus year
  (table in Section~\ref{subsec:keys}). Exception: Moen's undated
  preprint has key \texttt{Mo}, because its year is unknown.
\item \textbf{1.8.} Operators declared with \verb|\DeclareMathOperator|.
\item \textbf{1.9.} Every Draft~1 label is kept, except
  \texttt{eq:d}, which was a duplicate definition of $d$. The macro
  \verb|\ht| was renamed \verb|\hl|, since \verb|\ht| is a \TeX\
  primitive.
\item \textbf{2.} Order: abstract, introduction, literature review,
  background, notation, body, summary and open questions, glossary,
  acknowledgments, references. Draft~1 had the notation table before
  the introduction and no background section.
\item \textbf{2.3.} Literature review split at 1999. Works added:
  [Mc09], [Li12], [Lu20], [Sz24], [Ne86], [Si79].
\item \textbf{2.4, 3.3.} Each background subsection and each section
  has at least one example.
\item \textbf{3.2.} Every section and subsection opens with an
  orienting paragraph.
\item \textbf{3.8.} Prose has no semicolons. Sentences were shortened.
\item \textbf{5.1--5.5.} Abstract opens with the AI disclosure. Email
  added to the author line. Acknowledgments name each AI and its role,
  say that AI review is not certification, and point to this log.
\item \textbf{6.2.} Title restored to the 1999 title. Draft~1 had
  ``On the unitary dual of the metaplectic cover of $\mathrm{SL}(2,F)$,
  $F$ $p$-adic''. Subtitle ``Draft 2. Corrected and expanded version''.
  Date line ``Original draft: 1999. Corrected and expanded draft:
  2026.''
\item \textbf{6.3.} Note on the original manuscript in a
  \verb|\thanks| footnote.
\item \textbf{6.4.} Revision history removed from the body (the
  Draft~1 acknowledgments paragraph on revision history, Draft~1
  Remark~4.6, and the phrase ``the earlier draft'' in Section~6.3 of
  Draft~1).
\item \textbf{6.6.} Statements not fully verified are labeled
  Remark~2026.1--2026.5.
\end{itemize}

\paragraph{[DECIDE] items, defaults used.} Document class: kept
\texttt{article} from Draft~1. Times package: \texttt{newtx} with
fallback. Glossary at the end, before the acknowledgments. Author
line: Draft~1 listed Argon-Gemini and Sol-ChatGPT after ``with''.
Claude was added to that list because it supplied new proofs. Remove
it if you prefer.

\subsection{Mathematical changes and corrections}
\begin{enumerate}[leftmargin=*]
\item \textbf{Example 5.7 (double cover).} Draft~1 said ``This does not
  imply that the double cover splits over $A$.'' That is true as a
  statement of logic, but it suggests that the cover does not split
  over $A$. In the tame case it does split as an abstract group:
  $\varpi^ku\mapsto\hl(\varpi)^k\hl(u)$ is a homomorphism when $n=2$,
  since $\At$ is abelian and the cocycle is trivial on units. The
  splitting is not canonical. Replaced by the true statement that
  $t\mapsto\hl(t)$ is not a homomorphism.
\item \textbf{Theorem 8.3 hypothesis.} Draft~1 assumed $\gcd(p,n)=1$.
  Draft~2 assumes $\gcd(p,2n)=1$, because the proof uses the splitting
  of $K$ through Corollary~8.2. The formula itself is expected to hold
  for $p\nmid n$, but that is not proved here.
\item \textbf{Intertwining operator (Section 7.4).} Draft~1 wrote
  $T_s\varphi(g,a)$ in a model with an extra variable $a$ that was
  never defined. Replaced by the standard operator
  $\varphi\mapsto\int_F\varphi(\wt^{-1}(u(x),1)g)\,dx$, with the
  matrix $J_s$ described in words.
\item \textbf{Example 9.5 (degree three).} Draft~1 asserted the value
  of the constant term of $e_K$. Draft~2 derives it (Example~6.6) and
  checks the answer against Lemma~7.2 and the linear side
  (Examples~3.7, 9.3).
\item \textbf{Proposition 7.1.} Draft~1 also asserted temperedness.
  That is now in Remark~2026.2 as unverified for covers.
\item \textbf{Remark 2026.4.} Records exactly what the 1999 manuscript
  asserted about complementary series and the representations at
  $s=1/n$, and why it is not proved.
\end{enumerate}

\subsection{Removed material}\label{subsec:removed}
Draft~1 Remark~4.6 (``Correction of an earlier multiplicity claim'')
compared with the 1999 text, so it moved here. Its content: the 1999
lemma claimed a 2-1 map $\Pi(A_n)\to\Pi(\At)_{\mathrm{gen}}$ through
characters $\eta_j$ that differ on $C\setminus A_n$. For odd $n$,
$A_n=C$, so the two choices coincide and the claim fails. The correct
invariant is the genuine central character on $\ZA$.

\subsection{Bibliography corrections}
\begin{itemize}[leftmargin=*]
\item [Ma84a, Ma84b]: year 1987 corrected to 1984, and the pages of
  part~II corrected from 297--305 to 297--306.
\item [Mo88]: title corrected to ``Irreducibility of unitary principal
  series for covering groups of $SL(2,k)$''. Author C.~H.~Moen.
\item [Mo]: Draft~1 described it as a 1979 University of Chicago
  Ph.D. thesis. We could not confirm this, so the entry reverts to the
  1999 description, ``preprint''. Draft~1 also referred to Moen as
  ``her''. The pronoun was removed.
\item [Ta17]: author's name S.~Tang (Shiang Tang). The Draft~1
  summary was corrected. The paper computes Plancherel measures by a
  comparison method and constructs genuine central characters of the
  torus in the simply laced case.
\item [GS75]: title ``\dots for the metaplectic group'' (Draft~1 had
  ``on'').
\item [Jo99]: volume and pages 73 (1999) 332--340 confirmed. Older
  citations give ``273, 1--9'', which is wrong.
\end{itemize}

\subsection{Citation keys}\label{subsec:keys}
\begin{center}
\begin{tabular}{@{}ll@{\qquad}ll@{}}
\toprule
Draft 1 & Draft 2 & Draft 1 & Draft 2\\
\midrule
Ar & Ar80 & K & Ku67\\
BZ1, BZ2 & BZ76, BZ77 & KP & KP84\\
F & Fl80 & M1, M2 & Ma84a, Ma84b\\
FK & FK86 & Mo1, Mo2 & Mo, Mo88\\
G & Ge76 & Sch & Sc98\\
GS & GS75 & Szpruch19 & Sz19\\
GoldSz & GS16 & Tang & Ta17\\
J1, J2, J3, J4 & Jo88, Jo98, Jo99, Jo01 & (new) & BD84, Li12, Lu20,
Mc09,\\
 & & & Ne86, Si79, Sz24\\
\bottomrule
\end{tabular}
\end{center}

\section{Sources consulted and not consulted}

Rule~4.3. Level~2 means checked against a source actually consulted.
Here that source was usually an abstract, a publisher page or an arXiv
listing, found by web search on November 13, 2026. No full text was
read in this revision.

\begin{longtable}{@{}>{\raggedright\arraybackslash}p{0.15\linewidth}>{\raggedright\arraybackslash}p{0.80\linewidth}@{}}
\toprule
Key & What was checked\\
\midrule
\endhead
Ar80 & Level 2: bibliographic data and the introduction (irreducible
off $\operatorname{Re}s=0$ except at $s=\pm1/3$).\\
BD84 & Not consulted. The claim that the arguments extend to finite
central extensions comes from the 1999 manuscript.\\
BZ76, BZ77 & Not consulted. Precise theorem numbers for the geometric
lemma not given.\\
Fl80, FK86, KP84 & Bibliographic data only (FK86 and KP84 confirmed in
another paper's reference list).\\
Ge76, GS75 & Bibliographic data only. The endpoint $s=1/2$ statement in
Remark~2026.4 is from memory (Level 3).\\
GS16 & Level 2: journal data, the assumption $p\nmid n$, and
Theorem~5.2 (reducible genuine principal series from unitary data),
from the arXiv introduction.\\
Jo88 & Not re-checked. Volume 61 agrees with Draft~1 and with the
October 2026 corrected edition of Jo01. The 1999 manuscript has 66.\\
Jo98, Jo01, Jo99 & Level 2: journal data. Jo01's reference list
includes the 1999 manuscript as a preprint.\\
Ku67 & Not consulted. Formula (3.2) for the splitting of $K$ was
spot-checked by script (Remark~2026.1).\\
Li12 & Theorem numbers 2.6.4, 5.6.3, 5.8.10 are taken from the October
2026 review of Jo01, when the PDF was read. Not re-checked here.\\
Lu20 & Level 2: abstract.\\
Mc09 & Level 2: arXiv abstract and introduction (Stone--von Neumann
for the torus, Satake isomorphism). The published version (believed to
be in Progress in Math.~300) was not confirmed, so the arXiv version is
cited.\\
Mo & Not found. See above.\\
Mo88 & Level 2: publisher page and abstract ($p$ odd).\\
Ne86 & Not consulted. Standard properties of the Hilbert symbol.\\
Sc98 & Level 2: title page of the thesis.\\
Si79 & Not consulted. The 1999 manuscript cites p.~198 for the Weyl
integration formula. The page was not checked, so no page is given.\\
Sz19 & Level 2: journal data and abstract ($p\mid n$ allowed,
$4\nmid n$ assumed).\\
Sz24 & Level 2: arXiv abstract and introduction (reducible unitary
genuine principal series parametrized by nontrivial quadratic
characters, odd $n$, two summands).\\
Ta17 & Level 2: arXiv abstract.\\
\bottomrule
\end{longtable}

\section{Verification of the mathematics}\label{sec:levels}

Level~1 means an independent derivation by Claude in this revision.

\begin{itemize}[leftmargin=*]
\item Level 1: Lemmas 3.2, 5.1, 5.2, 5.4, 6.1, 6.3, 7.2, 8.1, 9.1,
  Theorems 5.6, 6.4, 8.3, 9.4, Corollary 8.2, Proposition 7.1(a), and
  all examples. The proofs in the paper are the derivations.
\item Level 1 modulo citations: Proposition 7.1(b) (Frobenius
  reciprocity and the geometric lemma for covers). Theorem 8.3 (Weyl
  integration formula for $SL(2,F)$). Theorem 6.4 Step~1 ($p$-adic
  inverse function theorem). Section 3.2 (formula (3.2) for the
  splitting of $K$).
\item Level 3 or cited only: Remarks 2026.2--2026.5.
\end{itemize}

\paragraph{Points a human referee should check first.}
\begin{enumerate}[leftmargin=*]
\item Formula (3.2) for the splitting of $K$, and that $\kappa=1$ on
  upper triangular matrices. Lemma~7.2 and Examples 6.6 and 9.5 depend
  on it.
\item The Haar normalization (3.4), and that the quotient measure on
  $\At\backslash\Gt$ is $du\,dk$.
\item Step~2 of the proof of Theorem~6.4.
\item The kernel argument in Lemma~8.1.
\end{enumerate}

\section{The verification script}\label{sec:script}

\texttt{scripts/kubota\_torus\_check.py} models $F=\mathbb Q_p$ with
$n\mid p-1$. It represents elements of $F^\times$ by rationals and uses
the tame formula (H5). All checks are exact. For
\[
(p,n)\in\{(7,3),(13,3),(5,4),(13,4),(7,6),(13,6),(3,2),(7,2)\}
\]
it checks:
\begin{enumerate}[leftmargin=*]
\item the cocycle identity for $\beta$ on random elements of
  $SL(2,\mathbb Q)$, which validates the implementation;
\item $\beta(h(a),h(b))=(b,a)_n$ (Lemma~3.2(a));
\item the commutator $(a,b)_n^{-2}$ (Lemma~5.1);
\item $\beta(w,w^{-1})=1$ and $\wt^{-1}\hl(t)\wt=\hl(t^{-1})$
  (Lemma~5.4(a), (b));
\item $\hl(t)^{-1}=(t,-1)_n\hl(t^{-1})$ (Lemma~5.4(c));
\item the center is $(F^\times)^d$ with index $d^2$ (Lemma~5.2);
\item $\nu$ is a homomorphism exactly when $n\mid d^2$ (Lemma~9.1 and
  Remark~9.6). It fails for $n=2$ and $n=6$, as expected;
\item $u\mapsto(\varpi,u)_n^2$ has order $d$ on units (Lemma~7.2);
\item the splitting identity for $\kappa$ on random elements of
  $SL(2,\mathbb Z_p)$ (Remark~2026.1).
\end{enumerate}
It also checks $(5,-1)_4=-1$ (Example~3.1). All checks pass, and the
script exits with status~0. Run time is about one second. The name ends
in \texttt{\_check.py}, so the repository's \texttt{run\_checks.py}
runs it with the other checks.

\section{Concordance of numbering}

Draft~1 numbers on the left, Draft~2 numbers on the right. Unlabeled
Draft~1 items are matched by content.

\begin{longtable}{@{}lll@{}}
\toprule
Item & Draft 1 & Draft 2\\
\midrule
\endhead
Sections: intro, literature & 1, 2 & 1, 2\\
Covering groups and conventions & 3 & 3.2--3.4 (Background)\\
Table of notation & unnumbered, before \S1 & 4\\
Covering torus, orbital integrals & 4, 5 & 5, 6\\
Parabolic induction & 6 & 7\\
Characters, transfer & 7, 8 & 8, 9\\
Conclusions, glossary & 9, 10 & 10, 11\\
Lemma radical & 4.1 & 5.2\\
Remark on centralizers & 4.2 & 5.3\\
Theorem torus & 4.3 & 5.6\\
Example 2 and 3 & 4.4 & 5.7\\
Example 4 and 6 & 4.5 & 5.8\\
Remark on multiplicity claim & 4.6 & removed (this log, \S\ref{subsec:removed})\\
Lemma vanishing & 5.1 & 6.1\\
Theorem orbital & 5.2 & 6.4\\
Remark Satake & 5.3 & 6.7\\
Proposition tempered & 6.1 & 7.1\\
Remark reducibility & 6.2 & Remark 2026.3\\
Theorem character & 7.1 & 8.3\\
Caution & 7.2 & 8.4\\
Corollary trace & 7.3 & 8.2\\
Lemma norm & 8.1 & 9.1\\
Definition matching & 8.2 & 9.2\\
Theorem transfer & 8.3 & 9.4\\
Example degree three & 8.4 & 9.5\\
Remark Weil & 8.5 & 9.6\\
Remark spherical & 8.6 & 9.7\\
Equations (3.1), (3.2) & kubota, genuine & (3.1), (3.3)\\
Equations (3.3), (3.4) & $D$, $\delta_B$ & (3.5), (3.6)\\
Equation (3.5) & $d$ & removed (defined in (5.2))\\
Equations (4.1)--(4.5) & & (5.1)--(5.5)\\
Equations (5.1)--(5.5) & OZ, Phi, index, Abel, surj & (6.1), (6.2),
(6.3), (6.5), (6.6)\\
Equations (6.1), (6.2) & & (7.1), (7.2)\\
Equations (7.1)--(7.4) & char, char-Z, trace & (8.3), (8.4), (8.1),
(8.2)\\
Equations (8.1)--(8.7) & & (9.1)--(9.7)\\
\bottomrule
\end{longtable}

\end{document}
